% ============================================================
%  Block C -- what survives a limit
% ============================================================
\rsection{Block C --- what survives a limit}

\begin{signpost}[The same questions, with a limit in front]
Let $f_n \to f$. Each property of Blocks A and B can now be asked twice: do
the $f_n$ have it, and does $f$? Pointwise convergence answers ``no'' almost
every time. Uniform convergence --- Rudin's 7.7 --- was invented precisely to
change the answers, and it changes two of the three.
\end{signpost}

\ritem{T13}{Definition}
\emph{$f_n \to f$ \textbf{pointwise} on $E$ if $f_n(x) \to f(x)$ for each
$x \in E$; \textbf{uniformly} if
$\sup_{x\in E}|f_n(x)-f(x)| \to 0$ --- one $N$ serving every $x$ at once.}

\begin{center}
\begin{tikzpicture}[x=1mm, y=1mm]
  \draw[axis] (0,0) -- (60,0);
  \draw[axis] (0,0) -- (0,34);
  \draw[guide] (0,26) -- (56,26);
  \node[mlbl, anchor=west] at (57,26) {$1$};
  % the epsilon tube about the limit function (0 on [0,1), 1 at 1)
  \fill[black!9] (0,0) rectangle (52,5);
  \node[lbl, anchor=west] at (0,-9)
    {shaded: the $\varepsilon$-tube about the limit function};
  \foreach \p/\lab in {1.4/{n=2}, 4/{n=5}, 9/{n=15}}{
    \draw[seg] plot[smooth, domain=0:52, samples=40]
      (\x, {26*(\x/52)^\p});}
  \node[mlbl, anchor=south east] at (50,20) {$x^{n}$};
  \node[ptfill, minimum size=2.2pt] at (52,26) {};
  \node[ptopen] at (52,0) {};
  \node[mlbl, anchor=north] at (52,-1.6) {$1$};
\end{tikzpicture}
\end{center}
\figcap{$f_n(x)=x^n$ on $[0,1]$ (E9). The pointwise limit is $0$ on $[0,1)$
and $1$ at $1$ --- discontinuous. However large $n$ is, the graph still
climbs out of any $\varepsilon$-tube about the limit before reaching $1$, so
the convergence is not uniform. Continuity is lost exactly because the
$N$ that works depends on $x$.}

\ritem{T14}{Theorem}
\emph{(Rudin 7.12.) If each $f_n$ is continuous and $f_n \to f$ uniformly,
then $f$ is continuous.}

\begin{proof}
Fix $x_0$ and $\varepsilon>0$. Choose $n$ with
$\sup|f_n-f| < \varepsilon/3$, then $\delta$ with
$|f_n(x)-f_n(x_0)| < \varepsilon/3$ for $|x-x_0|<\delta$. For such $x$,
\[ |f(x)-f(x_0)| \le |f(x)-f_n(x)| + |f_n(x)-f_n(x_0)| + |f_n(x_0)-f(x_0)|
   < \varepsilon . \]
\end{proof}

\noindent
The three terms are where the name ``$3\varepsilon$ argument'' comes
from.\kn{Notice exactly what uniformity buys: the \emph{first and third}
terms. Under pointwise convergence one can still make each of them small,
but only after fixing the point --- and the middle term then needs a $\delta$
depending on the $n$ that was chosen for that point, so the estimate never
closes. E9 is what that failure looks like.}

\ritem{T15}{Theorem}
\emph{(Rudin 7.16.) If each $f_n$ is integrable on $[a,b]$ and
$f_n \to f$ uniformly, then $f$ is integrable and}
\[ \int_a^b f = \lim_{n\to\infty}\int_a^b f_n . \]

\begin{proof}
Put $\varepsilon_n = \sup_{[a,b]}|f_n-f| \to 0$. Then
$f_n - \varepsilon_n \le f \le f_n + \varepsilon_n$, so for any partition
$U(P,f)-L(P,f) \le U(P,f_n)-L(P,f_n) + 2\varepsilon_n(b-a)$. Choosing $n$
with $2\varepsilon_n(b-a)$ small and then $P$ making the first term small
gives integrability by 6.6. The same inequalities bound
$\bigl|\int f - \int f_n\bigr|$ by $\varepsilon_n(b-a) \to 0$.
\end{proof}

\ritem{T16}{Theorem}
\emph{Uniform convergence does \textbf{not} preserve differentiability, and
even when the limit is differentiable, $\lim f_n'$ need not be $f'$.}

\begin{proof}
Take $f_n(x) = \sin(nx)/\sqrt n$ on $\R$ (E10). Then
$\sup|f_n| = 1/\sqrt n \to 0$, so $f_n \to 0$ uniformly and the limit is as
differentiable as anything could be. But
$f_n'(x) = \sqrt n\,\cos(nx)$, which is unbounded and has no limit at any
point --- at $x=0$, $f_n'(0)=\sqrt n \to \infty$ while $f' \equiv 0$.
\end{proof}

\begin{center}
\begin{tikzpicture}[x=1mm, y=1mm]
  \draw[axis] (0,14) -- (46,14);
  \foreach \p/\amp in {2/9, 5/5.7, 12/3.7}{
    \draw[seg] plot[smooth, domain=0:44, samples=120]
      (\x, {14 + \amp*sin(\p*\x*4)});}
  \node[lbl, anchor=north] at (22,-4) {$f_n = \sin(nx)/\sqrt n \to 0$};
  \begin{scope}[xshift=58mm]
  \draw[axis] (0,14) -- (46,14);
  \foreach \p/\amp in {2/5, 5/8, 12/12}{
    \draw[seg] plot[smooth, domain=0:44, samples=120]
      (\x, {14 + \amp*cos(\p*\x*4)});}
  \node[lbl, anchor=north] at (22,-4) {$f_n' = \sqrt n\,\cos(nx)$ diverges};
  \end{scope}
\end{tikzpicture}
\end{center}
\figcap{Why uniform convergence cannot control derivatives. On the left the
amplitude collapses like $1/\sqrt n$; on the right the same functions
differentiated, where the extra factor $n$ beats it and the amplitude grows
like $\sqrt n$. Uniform smallness of a function says nothing about its
slope --- a wiggle can be short and steep at the same time.}

\ritem{T17}{Theorem}
\emph{(Rudin 7.17 --- the repair.) Suppose each $f_n$ is differentiable on
$[a,b]$, the sequence $\bigl(f_n(x_0)\bigr)$ converges for some
$x_0 \in [a,b]$, and $(f_n')$ converges uniformly. Then $(f_n)$ converges
uniformly to some $f$, and $f' = \lim_n f_n'$.}

\noindent
The hypothesis has moved: it is the \emph{derivatives} that must converge
uniformly, and the functions need only converge at a single point.\kn{Read
T16 and T17 together as one lesson. Differentiation is not continuous with
respect to uniform convergence --- so if you want to differentiate a limit,
you must assume something about the derivatives, not about the functions.
The same asymmetry returns for series at 7.17's corollary, and it is why
term-by-term differentiation of a power series (8.1) needs its own theorem
while term-by-term integration does not.}

\ritem{T18}{Theorem}
\emph{Pointwise convergence preserves none of the three: not continuity, not
integrability, and not the value of the integral.}

\begin{proof}
Continuity fails by E9. For integrability, enumerate
$\Q \cap [0,1] = \{q_1,q_2,\dots\}$ and let $f_n$ be the indicator of
$\{q_1,\dots,q_n\}$; each $f_n$ is integrable, having finitely many
discontinuities (T10), and $f_n \to$ Dirichlet (E7) pointwise, which is not
integrable. For the integral, take the spikes of E11: $f_n \to 0$ pointwise
while $\int_0^1 f_n = 1$ for every $n$.
\end{proof}

\begin{center}
\begin{tikzpicture}[x=1mm, y=1mm]
  \draw[axis] (0,0) -- (56,0);
  \draw[axis] (0,0) -- (0,32);
  \foreach \w/\h/\lab in {24/8/{n=1}, 12/16/{n=2}, 6/32/{n=4}}{
    \draw[thin] (0,\h) -- (\w,\h) -- (\w,0);
    \node[mlbl, anchor=west] at (\w+1,\h) {$\lab$};}
  \node[lbl, anchor=west] at (16,24) {each has area $1$};
  \node[lbl, anchor=west] at (16,19.5) {and tends to $0$ at every $x$};
\end{tikzpicture}
\end{center}
\figcap{The spikes $f_n = n\cdot\mathbf 1_{(0,1/n)}$ (E11). At every fixed
$x>0$ the spike eventually moves left of $x$, so $f_n(x) \to 0$; at $x=0$ the
value is $0$ throughout. Yet $\int f_n = 1$ always. The mass does not
disappear --- it escapes by getting narrower and taller, which no pointwise
statement can see.}

\begin{deeper}[Two remarks the appendices already prepared]
\textbf{Pointwise limits are not arbitrary.} It is tempting to conclude that
pointwise convergence preserves nothing whatever. It preserves something:
a pointwise limit of continuous functions is of Baire class $1$, and not
every function is. Dirichlet's function is \emph{not} a pointwise limit of
continuous functions --- proved in Chapter 2's appendix at S2.13, using
Baire's theorem. So E7 arises as a pointwise limit of the \emph{integrable}
$f_n$ above, but never of continuous ones.

\smallskip
\textbf{And uniform convergence is not the only repair.} T15 requires
uniformity, but the conclusion $\int \lim = \lim \int$ survives under far
weaker hypotheses once the Lebesgue integral is available --- dominated
convergence needs only a pointwise limit and an integrable bound, which the
spikes of E11 conspicuously lack. Rudin reaches that at 11.32. The
uniformity in T15 is a limitation of the Riemann integral, not of the
statement.
\end{deeper}
